\chapter{Why Financial Machine Learning Is Different}\label{ch:ml:why-financial-machine-learning-is-different}

A new model forecasts next month's return of 500 stocks. Out of sample, over ten years, its R-squared is 0.57\%, and
the desk is pleased. It should be: in the simulated market the model was trained on, where the expected returns are
known because they were planted, the best forecast that can exist scores 0.60\%. The same gradient-boosted trees,
given a task whose signal is strong, explain 92\% of the variance. Nothing is wrong with either number. The model
did as well on the stocks as it did on the easy task; the stocks simply contain almost nothing to learn, and what they
contain changes, is shared by every stock at once, and is traded away by whoever learns it first. This chapter fixes
the vocabulary of machine learning the rest of the book uses, and measures, on data whose truth is known, the four
ways market data differ from the data the field grew up on.

\section{Learning from data: the vocabulary}

\begin{definition}[Supervised and unsupervised learning]\label{def:ml:why-financial-machine-learning-is-different:learning}
\emph{Supervised learning}\index{supervised learning} fits a function $f_\theta$ from features $x\in\mathbb R^p$ to
a target $y$ on examples $(x_i, y_i)$, $i = 1,\dots,n$, by minimising an empirical loss
$\mathcal L_n(\theta) = \frac1n\sum_i\ell(y_i, f_\theta(x_i))$, so that $f_\theta$ predicts $y$ for new $x$.
\emph{Unsupervised learning}\index{unsupervised learning} has no target: it describes the distribution of $x$ itself
(clusters, factors, a lower-dimensional representation, a density).
\end{definition}

The target is Book~7's prediction target (chapter~6): a forward return over a forecast horizon, or a label built from
one (chapter~2 of this book). The loss $\ell(y,\hat y)$ always takes two
arguments here; $\mathcal L_n$ is always subscripted. The \emph{learning rate} of later chapters is written $\eta$,
a symbol declared local to this book (the series keeps $\eta$ for the vol-of-vol elsewhere).

\begin{definition}[Training, validation and test sets; hyperparameter]\label{def:ml:why-financial-machine-learning-is-different:sets}
The \emph{training set}\index{training set} is the data on which $\theta$ is fitted. A
\emph{hyperparameter}\index{hyperparameter} is a choice the fit does not make itself (a penalty, a tree depth, a number
of layers, a learning rate); the \emph{validation set}\index{validation set} is data held out from the fit on which
hyperparameters are chosen. The \emph{test set}\index{test set} is data used once, after every choice is made, to
estimate how the chosen model will perform.
\end{definition}

On market data the three sets are consecutive stretches of time, in that order (\cref{fig:ml:why:split}), because the
model will be used on a future it has not seen. The \omterm{def:ml:why-financial-machine-learning-is-different:sets}{test set} is Book~7's holdout set (chapter~1): touched once. A \omterm{def:ml:why-financial-machine-learning-is-different:sets}{test
set} consulted twice has become a \omterm{def:ml:why-financial-machine-learning-is-different:sets}{validation set}.

\begin{omfigure}
\begin{tikzpicture}[x=0.95cm,y=1cm,font=\small]
\draw[->,thick] (0,0) -- (13.2,0) node[right]{time};
\fill[omDef!25] (0,0.2) rectangle (7,0.9);
\fill[omProp!30] (7.15,0.2) rectangle (9.5,0.9);
\fill[omThm!25] (9.65,0.2) rectangle (12.8,0.9);
\node at (3.5,0.55) {training: fit $\theta$};
\node at (8.33,0.55) {validation};
\node at (11.22,0.55) {test: used once};
\draw[omDef,thick] (3.5,-0.15) -- (3.5,-0.45) node[below,align=center]{\footnotesize features and targets\\[-2pt]\footnotesize known at each date};
\draw[omProp,thick] (8.33,-0.15) -- (8.33,-0.45) node[below,align=center]{\footnotesize choose the\\[-2pt]\footnotesize hyperparameters};
\draw[omThm,thick] (11.22,-0.15) -- (11.22,-0.45) node[below,align=center]{\footnotesize estimate the\\[-2pt]\footnotesize generalisation error};
\end{tikzpicture}

\omcaption{The three sets on market data are consecutive stretches of time. Chapter~3 replaces
the single split by purged folds and walk-forward schemes; the order never changes.}\label{fig:ml:why:split}
\end{omfigure}

\begin{definition}[Generalisation error, overfitting, bias--variance decomposition]\label{def:ml:why-financial-machine-learning-is-different:generalisation}
The \emph{generalisation error}\index{generalisation error} of a fitted model $\hat f$ is its expected loss on a new
draw $(x, y)$ from the distribution the model will meet. \emph{Overfitting}\index{overfitting} is the fitting of
features of the training sample that do not recur, so that the in-sample loss understates the generalisation error.
For the squared loss, the \emph{bias--variance decomposition}\index{bias--variance decomposition} splits the expected
generalisation error at a point into noise, squared bias and variance (\cref{prop:ml:why:bv}).
\end{definition}

\begin{proposition}[Bias--variance decomposition]\label{prop:ml:why:bv}
Let $y = \mu(x) + \varepsilon$ with $\E[\varepsilon\mid x] = 0$ and $\Var(\varepsilon\mid x) = \sigma^2(x)$, and let
$\hat f$ be fitted on a training sample independent of $(x, y)$. Then
\[
\E\bigl[(y - \hat f(x))^2 \bigm| x\bigr] = \sigma^2(x) + \bigl(\E[\hat f(x)] - \mu(x)\bigr)^2 + \Var\bigl(\hat f(x)\bigr),
\]
where $\E$ and $\Var$ of $\hat f(x)$ are over training samples.
\end{proposition}

\begin{proof}
Write $y - \hat f = \varepsilon + (\mu - \E\hat f) + (\E\hat f - \hat f)$. The three terms are uncorrelated given $x$:
$\varepsilon$ is independent of the training sample and has mean zero, and $\E\hat f - \hat f$ has mean zero over
training samples while $\mu - \E\hat f$ is a constant. Square and take expectations.
\end{proof}

The noise $\sigma^2$ is the floor no model goes below. In most of machine learning the noise is small and the fight is
between bias (a model too rigid for $\mu$) and variance (a model that follows the sample). On market data $\sigma^2$ is
almost all of the loss, and a model with any variance at all spends it on noise: the balance tips towards rigid,
regularised models (chapter~4).

\section{Signal-to-noise near zero and the R-squared ceiling}

\begin{definition}[Out-of-sample R-squared, predictability ceiling]\label{def:ml:why-financial-machine-learning-is-different:r2}
The \emph{out-of-sample R-squared}\index{out-of-sample R-squared} of forecasts $\hat y_i$ of returns $y_i$ on a \omterm{def:ml:why-financial-machine-learning-is-different:sets}{test
set} is
\[
R^2_{\mathrm{oos}} = 1 - \frac{\sum_i(y_i - \hat y_i)^2}{\sum_i y_i^2},
\]
measured against a forecast of zero, not against the sample mean. The \emph{predictability ceiling}\index{predictability ceiling}
of a dataset is $R^2_{\mathrm{oos}}$ of the true conditional expectation $\mu(x_i) = \E[y_i\mid x_i]$: the best any
model of those features can score.
\end{definition}

The zero benchmark follows Gu, Kelly and Xiu, whose study of 30\,000 stocks over 60 years uses it because the historical
mean of a stock's return is so noisy a forecast that beating it is too easy. Their monthly out-of-sample stock-level
R-squared runs from 0.16\% for a three-variable linear model to 0.33--0.40\% for trees and neural networks, with
0.40\% for the best network. Those are the numbers of a field where half a per cent is a strong result.

\begin{proposition}[The ceiling bounds every model; R-squared and correlation]\label{prop:ml:why:ceiling}
\begin{enumerate}
\item For any forecast $\hat f(x)$ built from the features, $\E(y - \hat f)^2 \ge \E(y - \mu)^2$: in expectation, no
model's $R^2_{\mathrm{oos}}$ exceeds the ceiling.
\item If $y$ and a forecast $f$ have mean zero, $\operatorname{Corr}(f, y) = \rho$, and the forecast is used as $c f$, then
$R^2_{\mathrm{oos}}(c) = (2c\rho\sigma_y\sigma_f - c^2\sigma_f^2)/\sigma_y^2$, maximal at $c^\star = \rho\sigma_y/\sigma_f$
where it equals $\rho^2$, and zero at $c = 2c^\star$.
\end{enumerate}
\end{proposition}

\begin{proof}
(1) $\E(y - \hat f)^2 = \E(y - \mu)^2 + \E(\mu - \hat f)^2$ because $y - \mu$ is uncorrelated with every function of $x$.
(2) Expand $\E(y - cf)^2 = \sigma_y^2 - 2c\rho\sigma_y\sigma_f + c^2\sigma_f^2$ and maximise the quadratic in $c$.
\end{proof}

The second part links R-squared to the information coefficient of Book~7 (chapter~6): a forecast whose correlation
with returns is $\rho$ earns at most $\rho^2$, and a correlation of 0.075 is an R-squared of 0.56\%. It also says how a
good forecast scores badly: a forecast with the right direction and twice the right size scores zero.

The book's sandbox, \texttt{firm.mlsynth}, draws a monthly cross-section of 500 stocks with twenty characteristics
published as cross-sectional ranks in $[-1, 1]$, as the empirical literature ranks them. Six carry signal, linearly and
through an interaction, a square, a threshold and an absolute value; fourteen are noise. Returns add a market factor,
ten industry factors and heavy-tailed specific shocks, and the planted expected return is scaled so that the ceiling
is 0.7\%. Ridge regression, LightGBM's gradient-boosted trees and a small multilayer perceptron are trained on twenty
years and tested on the next ten (\cref{tab:ml:why:models}); for comparison, the same three models are fitted to a task
with a Bayes-optimal R-squared of 0.95 (20\,000 examples, a smooth nonlinear function of five of twenty features).

\begin{table}[htbp]
\centering\small
\begin{tabular}{@{}lccccc@{}}
\toprule
 & \multicolumn{3}{c}{500 stocks, ceiling 0.60\% out of sample} & \multicolumn{2}{c}{strong signal, Bayes 0.95}\\
\cmidrule(lr){2-4}\cmidrule(l){5-6}
model & in sample & out of sample & corr.\ with the truth & in sample & out of sample\\
\midrule
ridge regression & 0.39\% & 0.38\% & 0.74 & 0.72 & 0.72\\
boosted trees & 1.45\% & 0.57\% & 0.92 & 0.92 & 0.92\\
neural network & 0.64\% & 0.28\% & 0.74 & 0.94 & 0.93\\
\bottomrule
\end{tabular}
\caption{Three models on two tasks: twenty years of training and ten of test on \texttt{firm.mlsynth}'s 500 stocks,
and 20\,000 training examples of a task with a strong signal. R-squared against zero for the stocks, against the mean for
the generic task. Data: \texttt{ml\_why.compare} and \texttt{ml\_why.strong\_signal}.}\label{tab:ml:why:models}
\end{table}

\Cref{tab:ml:why:models} is the chapter in miniature. On the strong task every model's in-sample and out-of-sample scores
agree and the flexible models win. On the stocks, the boosted trees' in-sample R-squared is 1.45\%, twice the ceiling
of the training years (0.72\%): the excess is noise it has fitted. Its out-of-sample 0.57\% is 95\% of the ceiling, the
best of the three; the linear model reaches 64\% of it, and the network, trained here without the care of
chapter~7, less than half. The forecasts' correlations with the truth
(0.74 to 0.92) are high; their correlations with the realised returns are below 0.08.

\begin{omfigure}
\begin{tikzpicture}
\begin{axis}[width=11cm,height=5.2cm,ybar,bar width=4pt,xlabel={monthly out-of-sample R-squared (\%)},ylabel={months},
  tick label style={font=\small},label style={font=\small},xmin=-1.1,xmax=2.85,ymin=0,grid=major,grid style={gray!25},
  legend columns=2,legend style={font=\footnotesize,at={(0.5,-0.32)},anchor=north,draw=none,fill=none,
  /tikz/every even column/.append style={column sep=8pt}},table/col sep=comma]
\addplot[fill=omDef!55,draw=omDef] table[x=x,y=boosting]{figdata/ml/01-why-financial-machine-learning-is-different/months.csv};
\addplot[fill=gray!45,draw=gray] table[x=x,y=truth]{figdata/ml/01-why-financial-machine-learning-is-different/months.csv};
\draw[omThm,thick,dashed] (axis cs:0,0) -- (axis cs:0,26);
\legend{boosted trees,the truth $\mu$}
\end{axis}
\end{tikzpicture}

\omcaption{The 120 test months, one R-squared per month over 500 stocks, for the boosted trees and for the planted truth
itself. Bars left of the dashed line are months where the forecast did worse than zero: 18\% of them for the model, and
some for the truth. Data: \texttt{ml\_why.compare} on \texttt{firm.mlsynth}.}\label{fig:ml:why:months}
\end{omfigure}

\section{Non-stationarity and adversarial feedback}

The distribution a model meets is not the one it was trained on. Some change is exogenous (a new regulation, a new
venue, a crisis). Some is caused by the model's own trade and its competitors': a predictor that is public, or that
several firms have found, is traded until its return is smaller than its cost. Book~7 measured the result on real
anomalies as post-publication decay (chapter~13) and its cause as crowding (chapter~28). In the physical sciences the
data do not read the paper.

\texttt{firm.mlsynth} can age a signal. From the first test month the strongest linear characteristic's effect decays
with a half-life of two years, and five years into the test the interaction between two characteristics changes sign.
The boosted trees, trained once on the first twenty years, keep their R-squared through the decay, which the other
characteristics hide, and lose it when the interaction flips (\cref{fig:ml:why:drift}): years 6 to 10 average
0.01\%, while the ceiling, which moves with the new truth, averages 0.83\%. Nothing in the model's inputs announced
the change; its errors did. Chapter~12 detects it and
chapter~27 builds the alarms.

\begin{omfigure}
\begin{tikzpicture}
\begin{axis}[width=11cm,height=5.4cm,xlabel={test year},ylabel={out-of-sample R-squared (\%)},
  tick label style={font=\small},label style={font=\small},xmin=0.5,xmax=10.5,ymin=-0.3,ymax=1.1,
  xtick={1,2,3,4,5,6,7,8,9,10},grid=major,grid style={gray!25},
  legend columns=3,legend style={font=\footnotesize,at={(0.5,-0.3)},anchor=north,draw=none,fill=none,
  /tikz/every even column/.append style={column sep=8pt}},table/col sep=comma]
\addplot[omDef,very thick,mark=*] table[x=year,y=stationary]{figdata/ml/01-why-financial-machine-learning-is-different/drift.csv};
\addplot[omThm,very thick,mark=square*] table[x=year,y=changing]{figdata/ml/01-why-financial-machine-learning-is-different/drift.csv};
\addplot[black!60,thick,dashed] table[x=year,y=ceiling_changing]{figdata/ml/01-why-financial-machine-learning-is-different/drift.csv};
\draw[gray,thin] (axis cs:0.5,0) -- (axis cs:10.5,0);
\legend{stationary market,signal decays then flips,ceiling (changing market)}
\end{axis}
\end{tikzpicture}

\omcaption{Boosted trees trained on years 1--20 of two panels with the same seed, scored year by year on the next ten.
In the second, characteristic~0's linear effect decays from test year~1 (half-life two years) and the interaction of
characteristics 0 and 3 flips sign at the start of year~6. Data: \texttt{ml\_why.drift\_by\_year}.}\label{fig:ml:why:drift}
\end{omfigure}

\section{Small effective samples}

The \omterm{def:ml:why-financial-machine-learning-is-different:sets}{training set} holds $240 \times 500 = 120\,000$ stock-months. It is not 120\,000 independent observations. The
stocks' unexpected returns share a market and an industry: their average pairwise correlation is 0.27, and the
effective number of independent stocks in a month, $n/(1 + (n-1)\bar\rho)$, is 3.7 out of 500. Whatever depends on the
market's path is estimated from 240 draws, not 120\,000.

The intercept is the clearest case. Trained on raw returns, a model's intercept is the training window's average market
return. The true premium in the sandbox is zero; the window happened to average $-0.56\%$ a month, about two standard
errors of a 240-month mean at a 4.5\% monthly volatility. The three models trained on raw returns score $-0.29\%$,
$-0.09\%$ and $-0.30\%$ out of sample, each of them worse than a forecast of zero. Trained on returns minus each month's
cross-sectional mean, they learn which stocks do better, not where the market goes, and score the table's numbers.

\begin{method}[Before fitting a model of returns]\label{met:ml:why:before}
\begin{enumerate}
\item Compute or bound the ceiling: on synthetic data exactly; on real data, from the best published R-squared or
information coefficient of comparable work (\cref{prop:ml:why:ceiling}).
\item Decide what the model is asked: the cross-section (demean or rank the target by date), or the level.
\item Count the effective sample: dates, not rows, for anything common to all assets; overlapping labels
(chapter~2) divide it again.
\item Fix the test period before looking at it, and plan how long a live track must be to tell the model from zero and
from the simpler model.
\end{enumerate}
\end{method}

Length is the other effective sample. The boosted trees' monthly R-squared has a mean of 0.62\% and a standard deviation
of 0.64\%: 4.3 months of live trading put its mean two standard errors from zero. Telling it from ridge regression takes
longer: the monthly difference has a mean of 0.21\% and a standard deviation of 0.47\%, which needs 21 months. With less
history the flexible model loses its advantage (\cref{fig:ml:why:learning}): trained on the last two years only, the
trees score 0.18\% and ridge 0.29\%; they cross between four and eight years. The more flexible the model, the more of
its training data goes to learning what the noise is not.

\begin{omfigure}
\begin{tikzpicture}
\begin{axis}[width=11cm,height=5.4cm,xlabel={training window (months)},ylabel={out-of-sample R-squared (\%)},
  tick label style={font=\small},label style={font=\small},xmin=0,xmax=252,ymin=0,ymax=0.7,xtick={24,48,96,144,240},
  grid=major,grid style={gray!25},legend columns=3,legend style={font=\footnotesize,at={(0.5,-0.3)},anchor=north,
  draw=none,fill=none,/tikz/every even column/.append style={column sep=8pt}},table/col sep=comma]
\addplot[omDef,very thick,mark=*] table[x=months,y=boosting]{figdata/ml/01-why-financial-machine-learning-is-different/learning.csv};
\addplot[omThm,very thick,mark=square*] table[x=months,y=ridge]{figdata/ml/01-why-financial-machine-learning-is-different/learning.csv};
\addplot[black!60,thick,dashed] table[x=months,y=ceiling]{figdata/ml/01-why-financial-machine-learning-is-different/learning.csv};
\legend{boosted trees,ridge regression,ceiling}
\end{axis}
\end{tikzpicture}

\omcaption{Learning curves on the stock panel: the same ten test years, training windows of the most recent 2 to 20
years. Data: \texttt{ml\_why.learning\_curve}.}\label{fig:ml:why:learning}
\end{omfigure}

\section{What carries over from the rest of machine learning, and what does not}

The algorithms carry over: least squares and its penalties (Book~4, chapter~16), trees, networks, gradient methods
(Book~4, chapter~24). So does the discipline of held-out evaluation. What does not carry over is the default settings
and the intuitions trained on images and text: shuffled cross-validation (Book~7, chapter~20, and
chapter~3), a validation score read as a fact rather than a draw, \omterm{def:ml:why-financial-machine-learning-is-different:sets}{hyperparameters} tuned until the
\omterm{def:ml:why-financial-machine-learning-is-different:sets}{validation set} agrees, an accuracy of 55\% taken for a weak model when it may be close to the ceiling, and a model fitted
once and trusted for years. Every chapter of this book is one of these corrections, measured.

\section{Tutorial: the 0.57\% model}\label{tut:ml:why-financial-machine-learning-is-different:tut}

\begin{tutorial}
\textbf{Goal.} Measure the ceiling of a synthetic stock panel, compare three models in and out of sample, and repeat on
a task with a strong signal. \textbf{End state:} the table, \cref{fig:ml:why:months,fig:ml:why:drift,fig:ml:why:learning}.
\begin{tutsteps}
\item \textbf{The panel.} Characteristics are ranks of persistent latent processes; the truth is a linear and a
nonlinear part, scaled so that $R^2_{\mathrm{oos}}(r,\mu)$ hits the ceiling.
\omcode{code/firm/mlsynth/firm_mlsynth.py}{131}{163}{The synthetic stock panel with a planted expected return.}\label{lst:ml:why:panel}
\item \textbf{The target and the three models.} Train on returns minus each month's cross-sectional mean; score against
zero.
\omcode{code/ml/01-why-financial-machine-learning-is-different/python/ml_why.py}{75}{91}{In- and out-of-sample R-squared, and the ceiling.}\label{lst:ml:why:compare}
\item \textbf{Run} \texttt{compare()}, \texttt{compare(target='raw')}, \texttt{strong\_signal()},
\texttt{learning\_curve()}, \texttt{effective\_stocks()}, \texttt{drift\_by\_year()} and \texttt{fig\_why.py}.
\end{tutsteps}
\textbf{What to change next.} Lower the strong task's signal-to-noise ratio until the Bayes R-squared is 1\% and see
which model wins with 20\,000 examples (exercise~7); raise
\texttt{nonlin} to 0.9 and see whether ridge still reaches half the ceiling.
\end{tutorial}

\section{Build: synthetic learning tasks}\label{bld:ml:why-financial-machine-learning-is-different:mlsynth}

\begin{build}
\bfield{Purpose} Every model of this book is scored against the best predictor that exists for its data.
\bfield{Interface} \texttt{task(n, p, snr, kind, seed)} with its Bayes R-squared; \texttt{PanelConfig(n, months, k,
seed, ceiling, nonlin, premium, mkt\_vol, ind\_vol, spec\_vol, t\_df, decay\_from, decay\_half\_life, regime\_at)};
\texttt{panel(cfg) -> Panel} with \texttt{X} (months, stocks, characteristics), \texttt{r}, \texttt{mu},
\texttt{beta}, \texttt{industry}, \texttt{mkt}, and \texttt{Panel.flat(months)}; \texttt{r2\_oos},
\texttt{ceiling}, \texttt{months\_to\_detect}.
\bfield{Rules} Characteristics at month $t$ are known at its end and predict the return of month $t+1$; the truth is
never an input; seeded and deterministic.
\bfield{Acceptance tests} \texttt{code/firm/mlsynth/tests/}: the truth scores the Bayes R-squared on the generic task;
the ceiling is hit on average over seeds; the noise characteristics carry (almost) nothing; decay and flip change the
truth; the R-squared and detection formulas on small cases.
\bfield{Stretch} Characteristics observed with error, so that the ceiling from the features is below the truth's;
missing values and listings that start and end, as in \texttt{firm.synthmkt}.
\end{build}

\begin{omsources}
\item S.~Gu, B.~Kelly and D.~Xiu, ``Empirical asset pricing via machine learning'', \emph{Review of Financial Studies}
33(5), 2020 (NBER Working Paper 25398, 2018, revised 2019, Table~1).
\item S.~Geman, E.~Bienenstock and R.~Doursat, ``Neural networks and the bias/variance dilemma'', \emph{Neural
Computation} 4(1), 1992.
\item T.~Hastie, R.~Tibshirani and J.~Friedman, \emph{The Elements of Statistical Learning}, 2nd ed., Springer, 2009.
\item R.~Israel, B.~Kelly and T.~Moskowitz, ``Can machines `learn' finance?'', \emph{Journal of Investment Management},
2020 (SSRN 3624052).
\item M.~López de Prado, \emph{Advances in Financial Machine Learning}, Wiley, 2018.
\end{omsources}

\section{Exercises}

\begin{exercise}[$\star$]\label{exo:ml:why-financial-machine-learning-is-different:1}
A forecast of next month's returns has a correlation of 0.05 with them and is optimally scaled. What is its
\omterm{def:ml:why-financial-machine-learning-is-different:r2}{out-of-sample R-squared}? What if it is used at three times the optimal scale?
\end{exercise}

\begin{exercise}[$\star$]\label{exo:ml:why-financial-machine-learning-is-different:2}
The unexpected returns of 500 stocks have an average pairwise correlation of 0.27. What is the effective number of
independent stocks in a month? With 0.05?
\end{exercise}

\begin{exercise}[$\star$]\label{exo:ml:why-financial-machine-learning-is-different:3}
A model's monthly R-squared has a mean of 0.5\% and a standard deviation of 0.8\%. How many months of independent
results put the mean two standard errors from zero? Three?
\end{exercise}

\begin{exercise}[$\star\star$]\label{exo:ml:why-financial-machine-learning-is-different:4}
The boosted trees score 1.45\% in sample, where the ceiling is 0.72\%. How can a model beat the best possible predictor
in sample, and what does the gap say about its out-of-sample score?
\end{exercise}

\begin{exercise}[$\star\star$]\label{exo:ml:why-financial-machine-learning-is-different:5}
The training window's average market return was $-0.56\%$ a month, and the mean squared monthly stock return is
$\E r^2 = 0.0093$. What is the standard error of a 240-month mean at a monthly volatility
of 4.5\%, and how many points of R-squared does an intercept error of $-0.56\%$ cost?
\end{exercise}

\begin{exercise}[$\star\star$]\label{exo:ml:why-financial-machine-learning-is-different:6}
\emph{Find the flaw.} ``We trained our model on raw monthly returns and its \omterm{def:ml:why-financial-machine-learning-is-different:r2}{out-of-sample R-squared} is negative, so the
characteristics carry no information about returns.''
\end{exercise}

\begin{exercise}[$\star\star\star$]\label{exo:ml:why-financial-machine-learning-is-different:7}
\emph{Coding.} Run \texttt{ml\_why.strong\_signal} with a signal-to-noise ratio that gives a Bayes R-squared of 1\%
(\texttt{snr = 0.0101/0.9899}), 20\,000 training examples. Report the three models' \omterm{def:ml:why-financial-machine-learning-is-different:r2}{out-of-sample R-squared} and explain
the ranking.
\end{exercise}

\begin{exercise}[$\star\star\star$]\label{exo:ml:why-financial-machine-learning-is-different:8}
Prove part 2 of \cref{prop:ml:why:ceiling} and use it to explain why a model whose forecasts correlate at 0.92 with
the truth can still score below zero out of sample.
\end{exercise}

\section{Problem: The 0.57\% Model}

\begin{problem}[{Weekend problem --- a model near its ceiling}]\label{pb:ml:why-financial-machine-learning-is-different:1}
The chapter's panel: 500 stocks, twenty characteristics, twenty years of training and ten of test, and a planted truth.

\medskip\noindent\textbf{Part I --- The ceiling.}
\begin{enumerate}
\item What is the \omterm{def:ml:why-financial-machine-learning-is-different:r2}{predictability ceiling} in the training and in the test years, and why do they differ?
\item What correlation with returns does an R-squared of 0.60\% correspond to?
\item Why is the R-squared measured against zero rather than the sample mean?
\item What share of the months does the truth itself lose to a forecast of zero, from \cref{fig:ml:why:months}?
\end{enumerate}

\medskip\noindent\textbf{Part II --- The models.}
\begin{enumerate}[resume]
\item Give the three models' in- and \omterm{def:ml:why-financial-machine-learning-is-different:r2}{out-of-sample R-squared}.
\item What share of the ceiling does each reach out of sample?
\item Why does the ranking on the strong task differ from the ranking on the stocks?
\item Which term of the \omterm{def:ml:why-financial-machine-learning-is-different:generalisation}{bias--variance decomposition} explains the trees' in-sample 1.45\%?
\end{enumerate}

\medskip\noindent\textbf{Part III --- Effective samples.}
\begin{enumerate}[resume]
\item How many independent stocks is a month worth, and why?
\item What do the models score when trained on raw returns, and why?
\item How long a training window does boosting need to beat ridge regression?
\item How many months of live results tell the trees from zero, and from ridge?
\end{enumerate}

\medskip\noindent\textbf{Part IV --- The verdict.}
\begin{enumerate}[resume]
\item What happens to the trees when the interaction flips sign, and what did the inputs show?
\item State the \emph{named result}: the ceiling, the model's share of it, and the months needed to tell it from zero
and from ridge.
\item Would you trade the trees or ridge regression, with ten years of history? With two?
\item What should a monitoring rule watch, given part IV's first question?
\item How would crowding appear in the sandbox?
\item Why is the ceiling unknowable on real data, and what replaces it?
\item What does Gu, Kelly and Xiu's best monthly R-squared of 0.40\% suggest about real ceilings?
\item In one sentence: why is financial machine learning different?
\end{enumerate}
\end{problem}

\section{Interview questions}

\begin{interviewq}[$\star$ \iqroles{researcher, mle}]\label{iq:ml:why-financial-machine-learning-is-different:1}
Your model of monthly stock returns has an \omterm{def:ml:why-financial-machine-learning-is-different:r2}{out-of-sample R-squared} of 0.5\%. Is that good?
\end{interviewq}

\begin{interviewq}[$\star$ \iqroles{researcher, mle}]\label{iq:ml:why-financial-machine-learning-is-different:2}
State the \omterm{def:ml:why-financial-machine-learning-is-different:generalisation}{bias--variance decomposition}. Where do models of returns sit on the trade-off, and why?
\end{interviewq}

\begin{interviewq}[$\star\star$ \iqroles{researcher}]\label{iq:ml:why-financial-machine-learning-is-different:3}
You have ten million rows of daily data on 2\,000 stocks over 20 years. How many independent observations do you have?
\end{interviewq}

\begin{interviewq}[$\star\star$ \iqroles{researcher}]\label{iq:ml:why-financial-machine-learning-is-different:4}
Why do return-prediction studies measure R-squared against a forecast of zero instead of the historical mean?
\end{interviewq}

\begin{interviewq}[$\star\star$ \iqroles{researcher, trader}]\label{iq:ml:why-financial-machine-learning-is-different:5}
A model had an information coefficient of 0.06 in its backtest and 0.02 in its first six months live. List the
possible causes and how you would tell them apart.
\end{interviewq}

\begin{interviewq}[$\star\star\star$ \iqroles{researcher, mle}]\label{iq:ml:why-financial-machine-learning-is-different:6}
Show that a correctly signed forecast can have a negative \omterm{def:ml:why-financial-machine-learning-is-different:r2}{out-of-sample R-squared}, and say what you would do about it.
\end{interviewq}
